\documentclass[11pt]{article}
\usepackage[margin=1in]{geometry}
\usepackage[T1]{fontenc}
\usepackage{lmodern,microtype,amsmath,amssymb,amsthm,booktabs}
\usepackage[hidelinks]{hyperref}
\hypersetup{pdftitle={Report 142 All orders for the weighted Dyck Newton diagonal},pdfauthor={},pdfsubject={A292692 fixed order asymptotics on the full interior ratio range}}
\pdfinfoomitdate=1\pdftrailerid{}\pdfsuppressptexinfo=15
\setlength{\emergencystretch}{2em}
\allowdisplaybreaks[2]
\newtheorem{theorem}{Theorem}[section]
\newtheorem{lemma}[theorem]{Lemma}
\newtheorem{proposition}[theorem]{Proposition}
\newtheorem{corollary}[theorem]{Corollary}
\theoremstyle{remark}\newtheorem{remark}[theorem]{Remark}
\newcommand{\LL}{\mathcal L}
\newcommand{\fall}[2]{#1^{\underline{#2}}}
\newcommand{\Q}{\mathbb Q}
\newcommand{\Z}{\mathbb Z}
\newcommand{\E}{\mathbb E}
\newcommand{\Prob}{\mathbb P}
\newcommand{\eps}{\varepsilon}
\newcommand{\dd}{\mathcal D}
\newcommand{\coef}[2]{[#1]#2}
\title{All orders for the weighted Dyck Newton diagonal}
\author{Report 142}
\date{October 2, 2026}
\begin{document}
\maketitle
\begin{abstract}
We prove a complete fixed-order asymptotic expansion for the Newton
coefficients of the weighted Dyck polynomials underlying OEIS A292692.
The moving ratio $\alpha=m/N$ may lie in any compact subset of $(0,1)$.
Every normalized coefficient has the form $\sqrt{1-s^2}$ times a rational
function of the saddle parameter $s$, where $\alpha=2s^2/(1+s)$.
An explicit finite algorithm determines the coefficients. The full
interior range follows by freezing a convolution inverse at the base
state, controlling the change of saddle at neighboring states, and
bootstrapping a bounded error through finitely many nested compact sets.
On the diagonal, all relative coefficients belong to $\Q(\sqrt{17})$;
in particular the second is
$-(30800095+3912441\sqrt{17})/80494592$.
We also obtain an all-orders Lambert $W$ and Newton inversion algorithm,
with the real uncertainty retained inside the integer ceilings.
Report 141 supplies the exact weighted-path foundation; all analytic
arguments required for this continuation are given here. The conclusions
are fixed-order statements, with no effective remainder constants or
claims at the two ratio boundaries.
\end{abstract}

\section{Definitions and main results}
For a Dyck path of semilength $N$, weight each peak at coordinates $(a,b)$
by $k+a/b$, and let $p_N(k)$ be the total path weight. The exact bridge
proved in Report 141~\cite{foundation} identifies these polynomials by
\begin{equation}\label{eq:recurrence}
 p_0(k)=1,\qquad
 p_N(k)=(k+2N-2)p_{N-1}(k)
       +\sum_{i=0}^{N-1}p_i(k)p_{N-1-i}(k).
\end{equation}
Our analytic starting point may equivalently be this recurrence. Write
\begin{gather}
 \LL_m P=[\fall{k}{m}]P(k)=\frac{\Delta^mP(0)}{m!},\qquad
 \gamma_N=4^{-N}\binom{2N}{N},\label{eq:Newton}\\
 z_N(k)=\gamma_N\prod_{i=0}^{N-1}(k+2i+1),\qquad
 Q_{N,m}=\LL_mz_N,\qquad
 R_{N,m}=\frac{\LL_mp_N}{2NQ_{N,m}}.\label{eq:normalization}
\end{gather}
The last definition applies for $N\ge1$ and $0\le m\le N$.
The sequence of interest is
\begin{equation}\label{eq:diagonaldef}
 a_n=\LL_np_{2n},\qquad a_0=1.
\end{equation}
The factor $1/m!$ in \eqref{eq:Newton} is essential.

Throughout the interior ratio range, set
\begin{equation}\label{eq:saddleparameters}
 \alpha=\frac mN=\frac{2s^2}{1+s},\quad 0<s<1,\quad
 t=1-s^2,\quad h=\frac{1-s}{s},\quad
 \mu=\frac1\alpha=\frac{1+s}{2s^2},\quad
 B=\frac{t(s+2)}{4s^4}.
\end{equation}
The map $s\mapsto\alpha$ is strictly increasing from $(0,1)$ onto $(0,1)$.
When convenient, a function of $s$ is also regarded as a function of
$\alpha$ through this map.

\begin{theorem}[All fixed orders on the full interior]\label{thm:allorders}
There are unique analytic coefficient functions $R_j$ on $0<s<1$, of the
form
\begin{equation}\label{eq:rationalform}
 R_j(s)=\sqrt{1-s^2}\,r_j(s),\qquad r_j\in\Q(s),\qquad r_0=1,
\end{equation}
regular throughout this interval, with the following property. For every
fixed integer $K\ge1$ and every compact $J\subset(0,1)$,
\begin{equation}\label{eq:allorders}
 R_{N,m}=\sum_{j=0}^{K-1}R_j(s(m/N))N^{-j}
                  +O_{J,K}(N^{-K})
\end{equation}
uniformly over integer pairs with $m/N\in J$. The finite algorithm in
Section~\ref{sec:algorithm} determines all coefficients exactly.
\end{theorem}

\begin{corollary}[The diagonal coefficient field]\label{cor:diagonal}
For every fixed $K\ge1$,
\begin{equation}\label{eq:diagonal}
 a_n=C d^n\frac{n!}{\sqrt n}
 \left(\sum_{j=0}^{K-1}b_jn^{-j}+O_K(n^{-K})\right),
\end{equation}
where
\begin{gather}
 d=\frac{51\sqrt{17}-107}{4},\qquad
 C=\frac{\sqrt{2+2/\sqrt{17}}}{\pi^{3/2}},\qquad b_0=1,\label{eq:constants}\\
 b_j\in\Q(\sqrt{17})\quad(j\ge0),\qquad
 b_1=-\frac{989}{2176}-\frac{907\sqrt{17}}{36992},\label{eq:b1}\\
 b_2=-\frac{30800095+3912441\sqrt{17}}{80494592}.
 \label{eq:b2}
\end{gather}
\end{corollary}
The first relative correction and the constants were established in
Report 141. The continuation here includes all fixed orders, the entire
interior ratio range, the exact coefficient field, and an integer-safe
inverse at every fixed order. No convergence of the expansion, uniformity
for $K$ growing with $N$, boundary result at $\alpha=0$ or $1$, effective
onset or error constant, or priority claim is made. Finite exact checks
in the companion corroborate algebra; the analytic remainder is proved
below.

\section{Exact transforms and the positive convolution}\label{sec:exact}
Report 141 proves the formal convolution identity
\begin{equation}\label{eq:convolution}
 2Nz_N=\sum_{j=1}^N p_jz_{N-j}.
\end{equation}
All ordinary coefficients of $p_N,z_N$ are nonnegative. Isolating the
$j=N$ term therefore gives $0\preceq p_N\preceq2Nz_N$ in ordinary
coefficientwise order. Since
$\LL_m k^d=\left\{\begin{smallmatrix}d\\m\end{smallmatrix}\right\}\ge0$,
Newton extraction preserves this order, and hence
\begin{equation}\label{eq:bounded}
 0\le R_{N,m}\le1.
\end{equation}
Here $Q_{N,m}>0$ throughout the valid range, as the next formula shows.

Define analytic functions on $|z|<1$ by
\begin{equation}\label{eq:fh}
 f(z)=(1-z)^{-1/2},\qquad h(z)=f(z)-1,\qquad
 C_{N,m}=[z^N]f(z)h(z)^m.
\end{equation}
The product definition gives
\[
 \sum_{N\ge0}\frac{z_N(k)}{(2N-1)!!}z^N
 =f(z)(1+h(z))^k
 =f(z)\sum_{m\ge0}\frac{\fall{k}{m}}{m!}h(z)^m.
\]
Thus, with $(-1)!!=1$,
\begin{equation}\label{eq:Q}
 Q_{N,m}=\frac{(2N-1)!!}{m!}C_{N,m}.
\end{equation}
All relevant coefficients are positive. The saddle condition
$t h'(t)/h(t)=N/m$ is exactly \eqref{eq:saddleparameters}.

For any polynomial $P$, Newton multiplication gives
\begin{equation}\label{eq:multiplication}
 \LL_m(z_\ell P)=\sum_{r=0}^{\ell}c_{\ell,r}(m)\LL_{m-r}P,
 \qquad c_{\ell,r}(m)=\frac{\Delta^r z_\ell(m-r)}{r!}.
\end{equation}
Indeed, expand $z_\ell(k)$ in the Newton basis centered at $q=m-r$,
and multiply each term by $\fall{k}{q}$; the identity
$\fall{k}{q}\fall{(k-q)}{r}=\fall{k}{q+r}$ proves the assertion.
Only admissible predecessor indices are included; otherwise the
transformed summand is zero and its kernel is defined to be zero.
For admissible predecessors the coefficients in \eqref{eq:multiplication}
are nonnegative.

Put $L=\lfloor N/2\rfloor$. Splitting \eqref{eq:convolution} into
disjoint halves yields the exact equation
\begin{align}
 R_{N,m}&=1-(\mathcal K_NR)_{N,m}-B_{N,m},\label{eq:causal}\\
 (\mathcal K_NR)_{N,m}
 &=\sum_{\ell=1}^{L}\sum_{r=0}^{\ell}
       a_{N,m,\ell,r}R_{N-\ell,m-r},\label{eq:exactkernel}\\
 a_{N,m,\ell,r}
 &=\left(1-\frac\ell N\right)c_{\ell,r}(m)
                 \frac{Q_{N-\ell,m-r}}{Q_{N,m}},\label{eq:kerneldef}\\
 B_{N,m}&=\sum_{j=1}^{N-L-1}
                   \frac{\LL_m(p_jz_{N-j})}{2NQ_{N,m}}.\label{eq:oppositeexact}
\end{align}
The possible even midpoint occurs once. Every principal predecessor has
size at least $N/2$.

\section{A uniform shifted saddle expansion}\label{sec:saddle}
Let $S$ be any compact subinterval of $(0,1)$ in the $s$ variable. Define
\begin{equation}\label{eq:cumulants}
 \dd=t\frac d{dt}=-\frac{1-s^2}{2s}\frac d{ds},\qquad
 \kappa_j=\dd^j\log h(t),\qquad g_j=\dd^j\log f(t).
\end{equation}
These are rational functions of $s$; in particular $\kappa_1=\mu$ and
$\kappa_2=B>0$.

\begin{lemma}[Shifted expansion with summable remainders]\label{lem:saddle}
Fix $A>0$ and $K\ge1$. Uniformly on $S$ and for
$0\le r\le\ell\le A\log m$,
\begin{align}
 C_{N-\ell,m-r}
 &=\frac{f(t)h(t)^mt^{-N}t^\ell h^{-r}}{\sqrt{2\pi mB}}
 \left(\sum_{j=0}^{K-1}D_j(\ell,r;s)m^{-j}
  +O\bigl(m^{-K}(1+\ell+r)^{M_K}\bigr)\right),\label{eq:shifted}
\end{align}
for some fixed exponent $M_K$. Here $D_0=1$ and
$D_j\in\Q(s)[\ell,r]$.
\end{lemma}
\begin{proof}
Set $\psi(\theta)=h(te^{i\theta})/h(t)$ and
$g(\theta)=f(te^{i\theta})/f(t)$. After extracting the factors outside the
parentheses in \eqref{eq:shifted}, the Fourier integral before its Gaussian
normalization is
\begin{equation}\label{eq:Fourier}
 \frac1{2\pi}\int_{-\pi}^{\pi}
 g(\theta)\psi(\theta)^{-r}e^{i\ell\theta}
 \exp\{m(\log\psi(\theta)-i\mu\theta)\}\,d\theta.
\end{equation}
Use the analytic logarithm only near zero. Elsewhere the exact integrand
is
\[
 g(\theta)\psi(\theta)^{m-r}e^{-i(N-\ell)\theta}.
\]
No global logarithm is needed. Moreover $h$ has no zero on $|z|=t>0$,
since $h(z)=0$ would imply $z=0$.

The function $\psi$ is the characteristic function of the integer law
$\Prob(X=j)=[z^j]h(z)t^j/h(t)$, $j\ge1$. Every positive degree has positive
mass, so the law has span one. Its exponential moments are uniform on $S$,
because $t$ stays in a compact subset of $(0,1)$. The mean and variance
are $\mu$ and $B$. Uniform aperiodicity gives $|\psi|\le q_*<1$ on any
fixed outer arc. On a sufficiently small fixed inner arc,
\[
 \Re(\log\psi(\theta)-i\mu\theta)\le-b\theta^2
\]
with $b>0$ uniform on $S$. Choose $0<\eta<1/6$. The outer arc is
exponentially small; the intermediate arc
$m^{-1/2+\eta}\le|\theta|\le\theta_0$ contributes
$O(\exp(-c m^{2\eta}))$. Polynomial logarithmic shifts do not change
these conclusions, since $m-r\ge m/2$ eventually.

On the central arc put $\theta=u/\sqrt m$ and $z=m^{-1/2}$. After
removing $e^{-Bu^2/2}$, the remaining analytic factor is
\begin{equation}\label{eq:Gaussianexp}
 \exp\left\{
 \sum_{j\ge1}\frac{g_j-r\kappa_j+\ell\mathbf1_{j=1}}{j!}(iu)^jz^j
 +\sum_{j\ge3}\frac{\kappa_j}{j!}(iu)^jz^{j-2}
 \right\}.
\end{equation}
Expand in the real variable $z$ through degree $2K-1$. Fixed derivatives
at zero are polynomials in $u,\ell,r$ over $\Q(s)$. Odd powers of $z$
give odd polynomials in $u$ and integrate to zero.

Here is a uniform Taylor remainder bound. For $0\le z'\le m^{-1/2}$,
$|u|\le m^\eta$, the saddle remainder exponent is
$O(z'|u|^3)=O(m^{-1/2+\eta})u^2$. Each required fixed derivative has
polynomial $u$ growth. The amplitude exponent has modulus bounded by
\[
 C(1+\ell+r)z'|u|+C(1+r)(z'u)^2.
\]
Young's inequality bounds its linear term by
$\delta u^2/2+C'(1+\ell+r)^2/m$, and the second summand is bounded
uniformly in the logarithmic range. The quadratic term is absorbed for
large $m$. Analytic derivatives are uniformly bounded in a fixed complex
neighborhood of the compact saddle set. Thus every derivative required
for the remainder is bounded by
\[
 C_K(1+\ell+r)^{M_K}(1+|u|)^{E_K}e^{\delta u^2},
 \qquad \delta<\tfrac14\min_S B.
\]
Taylor's integral remainder, multiplied by the Gaussian and integrated,
is $O(m^{-K}(1+\ell+r)^{M_K})$. Extending polynomial Gaussian moments
to the whole real line costs a superpolynomially small error. This proves
the asserted expansion and its uniform remainder.
\end{proof}

The proof is also a finite coefficient algorithm: expand
\eqref{eq:Gaussianexp} through $z^{2j}$ and apply
\begin{equation}\label{eq:Gaussianmoments}
 u^{2a}\longmapsto\frac{(2a-1)!!}{B^a},\qquad
 u^{2a+1}\longmapsto0.
\end{equation}
The coefficient of $z^{2j}$ is $D_j$. Only
$\kappa_1,\ldots,\kappa_{2j+2}$ and $g_1,\ldots,g_{2j}$ are needed.
For later use introduce the normalized ratio series
\begin{equation}\label{eq:V}
 V_{\ell,r}(x)=
 \frac{\sum_{j\ge0}D_j(\ell,r;s)x^j}
      {\sum_{j\ge0}D_j(0,0;s)x^j}.
\end{equation}
It has polynomial coefficients in $\ell,r$ over $\Q(s)$, satisfies
$V_{0,0}=1$ identically, and has the corresponding relative uniform
remainder when truncated at a fixed order.

\section{A positive estimate away from the endpoints}\label{sec:middle}
Define
\[
 U_{N,m,\ell}=\frac{\LL_m(z_\ell z_{N-\ell})}{Q_{N,m}}.
\]
\begin{lemma}[Uniform middle bound]\label{lem:middle}
On every compact $J\subset(0,1)$ there are constants $C<\infty$ and
$0<q<1$ such that
\begin{equation}\label{eq:middle}
 U_{N,m,\ell}\le CNq^{\min(\ell,N-\ell)}
 \quad(1\le\ell<N,\ m/N\in J).
\end{equation}
\end{lemma}
\begin{proof}
Choose an integer $H=m/(1-s)+O(1)$. On the compact ratio range,
$H>m$ eventually and $H/N=2s^2/(1-s^2)+O(N^{-1})$ stays in a compact
positive interval. Nonnegative Newton coefficients imply
\[
 \LL_m(z_\ell z_{N-\ell})\le
 \frac{z_\ell(H)z_{N-\ell}(H)}{\fall H m}.
\]
The evaluation must be at an integer: then higher falling factorials
are nonnegative or zero. By \eqref{eq:Q},
\[
 \frac{z_N(H)}{\fall H m\,Q_{N,m}}
 =\frac{[z^N]f(z)(1+h(z))^H}{\binom Hm C_{N,m}}.
\]
Positive-coefficient evaluation bounds the numerator by
$f(t)(1+h)^Ht^{-N}$. Lemma~\ref{lem:saddle}, with no shift and $K=1$,
bounds the denominator below by a constant times
$\binom Hm f(t)h^mt^{-N}/\sqrt m$. Put $p=h/(1+h)=1-s$. The ratio
is therefore at most
\[
 \frac{C\sqrt m}{\binom Hm p^m(1-p)^{H-m}}\le C'N.
\]
The last inequality follows from the uniform central binomial estimate:
$p$ stays away from $0,1$, $m-Hp=O(1)$, and Stirling's formula makes the
binomial probability at least $c/\sqrt H$.

For $\ell\le N/2$, the product formula gives
\[
 \frac{z_\ell(H)z_{N-\ell}(H)}{z_N(H)}
 =\frac{\gamma_\ell\gamma_{N-\ell}}{\gamma_N}
 \prod_{i=0}^{\ell-1}\frac{H+2i+1}{H+2(N-\ell+i)+1}.
\]
The gamma ratio is uniformly bounded: $\gamma_\ell\le1$ and
\[
 \log\frac{\gamma_{N-\ell}}{\gamma_N}
 \le\sum_{j=N-\ell+1}^N\frac1{2j-1}<\frac12.
\]
Each product factor is at most $(H+N)/(H+2N-1)$, which is bounded by
some $q<1$ on our compact range for all sufficiently large $N$. Symmetry
treats $\ell>N/2$, and a larger constant absorbs the finite remaining
sizes. This proves \eqref{eq:middle}.
\end{proof}
In \eqref{eq:exactkernel},
$\sum_r a_{N,m,\ell,r}=(1-\ell/N)U_{N,m,\ell}$.
Also $p_j\preceq2jz_j$ controls \eqref{eq:oppositeexact} by the same
bound. Consequently discarding either endpoint beyond $A\log N$ costs
at most $CN^2q^{A\log N}$. For any prescribed fixed power saving, choose
$A$ large enough. The sharp logarithmic endpoint estimates below, rather
than the extra factor $N$ in \eqref{eq:middle}, control the retained terms.

\section{Polynomial coefficients at both endpoints}\label{sec:endpoints}
\subsection{The Newton multiplication polynomial}
Let $d=\ell-r$, let $U$ be a sum of $r$ independent uniform variables
on $[0,1]$, and let $S_d$ be an independent uniformly chosen $d$-subset
of $\{0,\ldots,\ell-1\}$. Repeated integration of the finite difference
in \eqref{eq:multiplication} gives the exact identity
\begin{equation}\label{eq:Newtonintegral}
 c_{\ell,r}(m)=\gamma_\ell\binom\ell r m^d
 \E\prod_{i\in S_d}\left(1+\frac{2i+1-r+U}{m}\right).
\end{equation}
Each offset has absolute value at most $2\ell$. Truncating the expectation
after degree $K-1$ in $m^{-1}$ has absolute error at most
\begin{equation}\label{eq:productremainder}
 \frac{(2\ell^2/m)^K}{K!}\exp(2\ell^2/m).
\end{equation}
Indeed the sum of the absolute elementary products is bounded by the
corresponding exponential series.

Write this expectation as $U_{\ell,r}(x)=\sum_{q\ge0}u_q(\ell,r)x^q$.
For fixed $q$,
\begin{equation}\label{eq:u}
 u_q(\ell,r)=\frac{\fall{(\ell-r)}q}{\fall\ell q}\,
 \E e_q(1-r+U,3-r+U,\ldots,2\ell-1-r+U).
\end{equation}
Newton identities and sums of powers show that the elementary symmetric
expression is a polynomial in $\ell,r,U$ over $\Q$, of degree at most
$2q$. At $\ell=0,\ldots,q-1$ it vanishes identically in $r,U$.
It is therefore divisible by $\fall\ell q$. Moreover moments of $U$ are
polynomials in $r$, from
$\E e^{zU}=((e^z-1)/z)^r$. Thus $u_q$ is a genuine polynomial in
$\ell,r$. Polynomial continuation, together with $\fall{(\ell-r)}q$,
also gives the correct value when $\ell<q$. For example
\begin{equation}\label{eq:u12}
 u_0=1,\quad u_1=(\ell-r)(\ell-r/2),\quad
 u_2=\frac{(\ell-r)(\ell-r-1)(12\ell^2-12\ell r-4\ell+3r^2+r-4)}{24}.
\end{equation}

\subsection{Top ordinary coefficients and their boundary zeros}
Define $c_a(n)=[k^{n-a}]p_n(k)$ for $n\ge a$, and set $c_a(n)=0$
for $0\le n<a$. This notation is distinct from the two-index Newton
coefficient $c_{\ell,r}(m)$.
\begin{lemma}[Polynomial top coefficients]\label{lem:top}
For every $a\ge0$, $c_a(n)$ extends to a polynomial in $\Q[n]$ of degree
at most $2a$. For $a\ge1$ it is divisible by $\fall n a$. In particular,
\begin{equation}\label{eq:Ca}
 C_a(n):=\frac{c_a(n)}{\fall n a}\in\Q[n],\qquad
 \deg C_a\le a,\qquad C_0=1.
\end{equation}
The quotient is understood polynomially at its apparent zero denominators.
\end{lemma}
\begin{proof}
Here $c_0(n)=1$. For $a\ge1$, coefficient extraction from
\eqref{eq:recurrence} gives the constructive recurrence
\begin{align}
 c_a(n)-c_a(n-1)
 &=(2n-2)c_{a-1}(n-1)\notag\\
 &\quad+\sum_{i=0}^{n-1}\sum_{b=0}^{a-1}
                   c_b(i)c_{a-1-b}(n-1-i),\qquad c_a(0)=0.
 \label{eq:toprec}
\end{align}
Assume the conclusion for smaller $a$. The right side is a polynomial
of degree at most $2a-1$: summing a monomial
$i^v(n-1-i)^w$ has degree at most $v+w+1\le2a-1$.
Its discrete antiderivative with value zero at $0$ is therefore a polynomial
of degree at most $2a$.

For an integer $1\le n<a$, the first term in \eqref{eq:toprec} vanishes.
Both factors of a convolution summand can be nonzero only if
$i\ge b$ and $n-1-i\ge a-1-b$, which would require $n\ge a$.
Every increment at $1\le n<a$ is therefore zero. Together with $c_a(0)=0$
this proves zeros at $0,\ldots,a-1$, and hence divisibility by
$\fall n a$. The same coefficient extraction, including the zero
convention, shows that this polynomial agrees with the actual top
coefficients. The asserted degree of the quotient follows.
\end{proof}
For example $c_1(n)=(3n^2-n)/2$. The boundary zeros in this lemma are
needed for the opposite-endpoint expansion; a degree bound alone would
not justify the coefficient algorithm at small integer indices.

\subsection{The opposite endpoint to every fixed order}
For a monomial $k^d$ and $0\le r\le d$,
\begin{equation}\label{eq:W}
 \frac{\Delta^r k^d(m-r)}{r!}
 =\binom dr m^{d-r}W_{d,r}(m^{-1}),\qquad
 W_{d,r}(x)=\sum_{q\ge0}\binom{d-r}q\E(-r+U)^q x^q.
\end{equation}
At every fixed order its coefficients are polynomial in $d,r$.
For valid integer indices it is an exact finite sum, with a remainder
bound of the form \eqref{eq:productremainder}. Outside that range below,
only its polynomial continuation multiplied by a vanishing falling
factorial is used.

Put $\eps=N^{-1}$ and define formal series
\begin{equation}\label{eq:FG}
 F_\ell(\eps)=\prod_{i=0}^{\ell-1}(1-(i+1/2)\eps)^{-1},\qquad
 G_r(x)=\prod_{i=0}^{r-1}(1-ix).
\end{equation}
Each fixed coefficient is polynomial in the indices, by logarithmic
expansion and Faulhaber sums. Their fixed-order remainders, when the
indices are logarithmic in $N$, are a fixed polynomial in $1+\ell$
times the requested power of $N^{-1}$.
Define
\begin{align}
 \beta_{j,r}&=s^{2j}\binom jr s^r(1-s)^{j-r},\label{eq:beta}\\
 Z_{j,r}(\eps)&=\sum_{a\ge0}(\mu\eps)^a C_a(j)
                       \fall{(j-r)}a W_{j-a,r}(\mu\eps).
 \label{eq:Z}
\end{align}
Only finitely many $a$ contribute at a prescribed order. The identity
\[
 c_a(j)\binom{j-a}r=C_a(j)\fall{(j-r)}a\binom jr
\]
justifies \eqref{eq:Z} when $j\ge a$ and $r\le j-a$.
For $j<a$ or $r>j-a$, its right side is zero at the integer indices used
in the sum, so the boundary convention is exact.

The formal expansion of the opposite endpoint is
\begin{align}
 B(\eps,s)&=\frac\eps2\sum_{j\ge1}\sum_{r=0}^j
  \beta_{j,r}F_j(\eps)G_r(\mu\eps)V_{j,r}(\mu\eps)Z_{j,r}(\eps)
   =\sum_{q\ge1}B_q(s)\eps^q.\label{eq:Bseries}
\end{align}
Every $B_q$ is rational in $s$. At fixed order each summand is a polynomial
in $j,r$ over $\Q(s)$, and its moments follow by Euler differentiation of
\begin{equation}\label{eq:betasymbol}
 \sum_{j\ge1,r}\beta_{j,r}x^jy^r
       =\frac1{1-s^2x(1-s+sy)}-1.
\end{equation}
All moments converge locally uniformly for $0<s<1$. In particular
$B_1=s^2/(2t)$.

We now justify the expansion with a genuine uniform remainder. For
$j\le A\log N$, $0\le d\le j$, and sufficiently large $N$, the
nonnegative monomial Newton coefficient is at most
$\binom dr m^{d-r}$, since $0\le m-r+U\le m$ in the integral formula.
Lemma~\ref{lem:saddle} and the factorial ratios then give
\begin{align}
 T_{j,d}:=\frac{\LL_m(k^dz_{N-j})}{Q_{N,m}}
 &\le C(2N)^{-j}m^dt^j(1+h^{-1})^d
  =Cs^{2j}M^{d-j},\label{eq:monomialbound}\\
 M&=m(1+h^{-1})\ge cN.\notag
\end{align}
Coefficientwise positivity implies
\[
 0\le c_a(j)\le2j\gamma_j e_a(1,3,\ldots,2j-1)
               \le2j\gamma_j\frac{j^{2a}}{a!}.
\]
Hence losses $a\ge H$ contribute at most
\begin{equation}\label{eq:lossbound}
 \sum_{a=H}^j c_a(j)T_{j,j-a}
 \le Cs^{2j}j\frac{(Cj^2/N)^H}{H!}\exp(Cj^2/N)
 =O\bigl(s^{2j}(1+j)^{2H+1}N^{-H}\bigr).
\end{equation}
For the finitely many retained degree losses, apply
Lemma~\ref{lem:saddle}, \eqref{eq:W}, and the finite-product expansions.
Every error is bounded by the requisite power of $N^{-1}$ times a fixed
polynomial in $j,r$ and the positive weight $\beta_{j,r}$.
All polynomial moments of $s^{2j}$ are uniformly summable on a compact
$s$ interval, so summation loses no logarithm. The extra factor $1/(2N)$
in \eqref{eq:Bseries} is retained. Lemma~\ref{lem:middle} disposes of
$j>A\log N$. Thus for every fixed $K$,
\begin{equation}\label{eq:Bexpansion}
 B_{N,m}=\sum_{q=1}^{K-1}B_q(s)N^{-q}+O_{J,K}(N^{-K})
\end{equation}
on every compact interior ratio range.

\subsection{The principal endpoint}
Combining \eqref{eq:Newtonintegral}, \eqref{eq:V}, \eqref{eq:FG}, and
\eqref{eq:Q} gives, uniformly for $\ell\le A\log N$,
\begin{align}
 a_{N,m,\ell,r}
 &=a_s(\ell,r)\left(\sum_{q=0}^{K-1}A_q(\ell,r;s)N^{-q}
       +O\bigl(N^{-K}(1+\ell)^{M_K}\bigr)\right),\label{eq:Aexpansion}\\
 a_s(\ell,r)&=\gamma_\ell s^{2\ell}\binom\ell r s^r(1-s)^{\ell-r},
 \label{eq:leadingkernel}\\
 A_{\ell,r}(\eps,s)&=(1-\ell\eps)F_\ell(\eps)
 G_r(\mu\eps)U_{\ell,r}(\mu\eps)V_{\ell,r}(\mu\eps)
     =\sum_{q\ge0}A_q(\ell,r;s)\eps^q.\label{eq:Aseries}
\end{align}
All $A_q$ are in $\Q(s)[\ell,r]$ and $A_0=1$. The normalization includes
\begin{equation}\label{eq:zerocancellation}
 A_{0,0}(\eps,s)=1,\qquad A_q(0,0;s)=0\quad(q\ge1).
\end{equation}
The remainder in \eqref{eq:Aexpansion} is relative to the positive
$a_s$. Its polynomially weighted sums are uniformly finite, since
$\sum_r a_s(\ell,r)=\gamma_\ell s^{2\ell}$.

\section{The exact coefficient algorithm and rationality}\label{sec:algorithm}
Let $H(\eps,\alpha)=\sum_{q\ge0}R_q(\alpha)\eps^q$.
Insert the endpoint series into the formal equation
\begin{equation}\label{eq:formalH}
 H(\eps,\alpha)=1-B(\eps,\alpha)
 -\sum_{\ell\ge1,r}a_s(\ell,r)A_{\ell,r}(\eps,s)
 H\left(\frac\eps{1-\ell\eps},
          \frac{\alpha-r\eps}{1-\ell\eps}\right).
\end{equation}
Every coefficient uses finitely many polynomial moments. The leading mass
is
\begin{equation}\label{eq:mass}
 q_0(s)=\sum_{\ell\ge1,r}a_s(\ell,r)=t^{-1/2}-1.
\end{equation}
Thus $(1+q_0)R_0=1$, giving $R_0=\sqrt t$. At order $q\ge1$, the only
new unknown on the right is $-q_0R_q$. If $P_q^{\rm known}$ denotes
the known order-$q$ principal sum after omitting that term, then
\begin{equation}\label{eq:Rrec}
 R_q=-\sqrt t\bigl(B_q+P_q^{\rm known}\bigr).
\end{equation}
All other terms use only $R_0,\ldots,R_{q-1}$ and finitely many
$\alpha$ derivatives. The conversion is
\begin{equation}\label{eq:alphaderivative}
 \partial_\alpha=\frac{(1+s)^2}{2s(s+2)}\partial_s.
\end{equation}
This constructs unique analytic coefficient functions on the whole
interior: all series moments converge locally uniformly and all
necessary denominators $B,t,\alpha,1+q_0$ are positive there.

\begin{proposition}[The zero index removes the radical obstruction]
Every $R_q$ constructed by \eqref{eq:Rrec} belongs to $\sqrt t\,\Q(s)$.
\end{proposition}
\begin{proof}
Suppose this holds for smaller indices. Every $\alpha$ derivative of
$\sqrt t$ times a rational function has the same form. Divide the known
principal residual summand by $\sqrt t$; it becomes a polynomial
$P_q(\ell,r;s)\in\Q(s)[\ell,r]$. At $\ell=r=0$ there is no kernel
correction by \eqref{eq:zerocancellation}, no argument shift, and no
denominator shift. Because the unknown $R_q$ contribution has been
omitted, we have
\begin{equation}\label{eq:Pzero}
 P_q(0,0;s)=0.
\end{equation}
All moments including $\ell=0$ are Euler derivatives of
\[
 \mathcal G(x,y)=(1-s^2x(1-s+sy))^{-1/2}.
\]
Evaluated at $x=y=1$, each derivative is $t^{-1/2}$ times a rational
function of $s$. Equation~\eqref{eq:Pzero} means that removing $\ell=0$
changes nothing. The principal residual, including its factor $\sqrt t$,
is therefore rational in $s$. The opposite residual $B_q$ is rational
by \eqref{eq:betasymbol}. Equation~\eqref{eq:Rrec} proves the result.
Without \eqref{eq:Pzero}, the removal of the zero index would introduce
an extra radical term, so this cancellation is essential.
\end{proof}

For clarity, the complete finite algorithm at a requested order is:
\begin{enumerate}
\item Form the necessary cumulants and Gaussian coefficients using
\eqref{eq:cumulants}, \eqref{eq:Gaussianexp}, and \eqref{eq:Gaussianmoments};
then divide the two truncated series in \eqref{eq:V}.
\item Generate $U_{\ell,r}$ by Newton identities and uniform-sum moments,
and generate $F_\ell,G_r$ by their finite-product logarithms.
\item Compute the finitely many $C_a$ needed by the discrete
antiderivative recurrence \eqref{eq:toprec}, and form \eqref{eq:Z}.
\item Compute the required principal and opposite moments by Euler
derivatives of $\mathcal G$ and \eqref{eq:betasymbol}.
\item Extract successive coefficients of \eqref{eq:formalH}, dividing
only by $1+q_0=t^{-1/2}$ as in \eqref{eq:Rrec}.
\end{enumerate}
All these operations are exact rational-function algebra and finite
series operations. The coefficients are regular in the full interior;
any apparent rational pole there is removable by the analytic construction.

\section{The uniform remainder on the full interior}\label{sec:fullrange}
The principal mass $q_0$ need not be less than one. We instead use an
exact inverse for a frozen convolution, with an explicit correction for
neighboring saddle changes.

\subsection{The frozen inverse and its moments}
For a fixed $s\in(0,1)$ define operators on bounded functions on $\Z^2$:
\begin{align}
 (\mathcal A_s E)(N,m)
 &=\sum_{\ell\ge1}\sum_{r=0}^{\ell}a_s(\ell,r)E(N-\ell,m-r),
 \label{eq:Aoperator}\\
 (\mathcal B_s E)(N,m)
 &=\sum_{\ell\ge1}\sum_{r=0}^{\ell}b_s(\ell,r)E(N-\ell,m-r),\qquad
 b_s(\ell,r)=\frac{a_s(\ell,r)}{2\ell-1}.
 \label{eq:Boperator}
\end{align}
These $b_s$ are inverse-kernel weights, distinct from the diagonal
coefficients $b_j$. Their two-variable symbols satisfy
\begin{gather}
 1+A_s(u,v)=[1-s^2u(1-s+sv)]^{-1/2},\label{eq:Asymbol}\\
 B_s(u,v)=1-[1-s^2u(1-s+sv)]^{1/2},\label{eq:Bsymbol}\\
 (I-\mathcal B_s)(I+\mathcal A_s)=I.\label{eq:inverseoperator}
\end{gather}
The orientation and signs follow by multiplication of the displayed
symbols; shifts compose as $T_{\beta}T_{\eta}=T_{\beta+\eta}$. Both
kernels are nonnegative, and
\begin{equation}\label{eq:kernelmasses}
 \|\mathcal A_s\|_1=t^{-1/2}-1,\qquad
 \|\mathcal B_s\|_1=1-\sqrt t<1.
\end{equation}
For bounded $E$, the double sum in the product is dominated by
$\|E\|_\infty\|\mathcal A_s\|_1\|\mathcal B_s\|_1$.
Consequently \eqref{eq:inverseoperator} is an absolutely convergent
identity on the entire lattice, not merely a formal one. Its inverse
norm is at most $1+\|\mathcal B_s\|_1<2$.

On any compact $S\subset(0,1)$ every polynomial moment of both kernels
is uniformly finite. Their first $s$ derivatives also have exponentially
decreasing, polynomially weighted tails. Indeed
\begin{equation}\label{eq:kernelderivative}
 \partial_s\log a_s(\ell,r)
 =\partial_s\log b_s(\ell,r)
 =\frac{2\ell+r}{s}-\frac{\ell-r}{1-s},
\end{equation}
so $|\partial_s a_s|\le C_S\ell a_s$ and similarly for $b_s$.
Summing in $r$ leaves $\gamma_\ell s^{2\ell}$ or
$\gamma_\ell s^{2\ell}/(2\ell-1)$, both geometrically decreasing
uniformly on $S$. In particular,
\begin{equation}\label{eq:Lipschitz}
 \|\mathcal A_s-\mathcal A_{s'}\|_1\le C_S|s-s'|
\end{equation}
when the intervening interval lies in $S$.

\subsection{A bounded extension and its local residual}
Fix $K$. Choose a compact interval $J_*$ strictly inside $(0,1)$,
containing the desired ratio range and all the nested intervals to be
used below, with additional margin. Put
\[
 H_K(N,m)=\sum_{j=0}^{K-1}R_j(m/N)N^{-j}.
\]
Extend the finitely many coefficient functions smoothly and boundedly
to $[0,1]$, agreeing with their true values on a neighborhood of $J_*$.
At valid states set $E_{N,m}=R_{N,m}-H_K(N,m)$; at every other lattice
state, including $N\le0$, set $E=0$. Equation~\eqref{eq:bounded} shows
that this is a bounded function on $\Z^2$.

The local expansions, the coefficient cancellations in
\eqref{eq:formalH}, and Lemma~\ref{lem:middle} prove the uniform residual
\begin{equation}\label{eq:localresidual}
 E(x)+(\mathcal K_xE)(x)=\rho(x),\qquad
 |\rho(x)|\le C N^{-K},\qquad x=(N,m),\quad m/N\in J_*.
\end{equation}
Here $\mathcal K_x$ is the finite exact kernel
\eqref{eq:exactkernel}. To justify this subtraction independently of any
remainder theorem, use Taylor expansion only for jumps at most
$A\log N$, whose ratios remain within the extra compact margin. All
coefficient-function derivatives are uniformly bounded there. Their
remainders are a fixed polynomial in the jump times $N^{-K}$, summable
against $a_s$. The long exact jumps act on bounded extensions and are
$O(N^{-K})$ by \eqref{eq:middle}; the formal tails have the same property
by geometric decay. Equation~\eqref{eq:Bexpansion} supplies the opposite
residual. Inadmissible zero-kernel terms are omitted before evaluating
$H_K$. No equation outside the local ratio domain is asserted.

Let $L_N=\lfloor A\log N\rfloor$, with $A$ sufficiently large for all
fixed-order tail bounds. From \eqref{eq:Aexpansion},
\begin{equation}\label{eq:perturbation}
 |a_{N,m,\ell,r}-a_{s(x)}(\ell,r)|
 \le C N^{-1}(1+\ell)^D a_{s(x)}(\ell,r)
 \quad(\ell\le L_N).
\end{equation}
The exact and frozen tails acting on bounded $E$ are $O(N^{-K})$.
Writing $\mathcal D_x$ for $\mathcal A_{s(x)}-\mathcal K_x$ restricted
to these short shifts, we obtain
\begin{equation}\label{eq:naturalfrozen}
 ((I+\mathcal A_{s(x)})E)(x)
   =\rho(x)+(\mathcal D_xE)(x)+O(N^{-K}).
\end{equation}
The sign of this perturbation is fixed by adding
$\mathcal A_{s(x)}-\mathcal K_x$ to \eqref{eq:localresidual}.

\subsection{A one power improvement}
\begin{lemma}\label{lem:improvement}
Let $J\subset\operatorname{int}J'\subset(0,1)$ be compact intervals,
with a further margin for the residual estimates. If, for an integer
$p\ge0$, the same error $E=R-H_K$ satisfies
\[
 |E_{N,m}|\le C_pN^{-p}\qquad(m/N\in J'),
\]
then it satisfies
\begin{equation}\label{eq:improvement}
 |E_{N,m}|\le C'_pN^{-\min(K,p+1)}\qquad(m/N\in J).
\end{equation}
\end{lemma}
\begin{proof}
Fix $x=(N,m)$ in the smaller range and freeze $s_0=s(x)$ for the
entire following convolution. Define, at \emph{every} lattice state $y$,
\[
 F(y)=((I+\mathcal A_{s_0})E)(y).
\]
This function is globally bounded, uniformly for $x$ in the fixed compact
range. The exact identity \eqref{eq:inverseoperator} gives
\begin{equation}\label{eq:frozenatbase}
 E(x)=F(x)-\sum_{b\ge1}\sum_{c=0}^b
                       b_{s_0}(b,c)F(x-(b,c)).
\end{equation}
The inverse tail $b>L_N$ is $O(N^{-K})$ against this bounded function.
For a retained shift $\beta=(b,c)$, including $(0,0)$, set
$y=x-\beta$ and $s_y=s(y)$. Its size is $N_y=N-O(\log N)$ and
\begin{equation}\label{eq:drift}
 |\alpha_y-\alpha_x|
 =\frac{|\alpha_xb-c|}{N-b}\le C\frac bN,
 \qquad |s_y-s_0|\le C\frac bN.
\end{equation}
For sufficiently large $N$, $y$ and every state reached by a second short
shift remain in $J'$ and have size comparable to $N$.

The natural neighboring equation has parameter $s_y$, not $s_0$.
Retain the difference explicitly:
\begin{align}
 F(y)
 &=((I+\mathcal A_{s_y})E)(y)
                  +((\mathcal A_{s_0}-\mathcal A_{s_y})E)(y)\notag\\
 &=O(N^{-K})+(\mathcal D_yE)(y)
                  +((\mathcal A_{s_0}-\mathcal A_{s_y})E)(y).
 \label{eq:neighbor}
\end{align}
The first short perturbation is $O(N^{-p-1})$ by
\eqref{eq:perturbation}, the induction hypothesis on all short
predecessors, and the finite weighted moments of $a_{s_y}$.

For the parameter-difference term, first discard shifts $\ell>L_N$
in both frozen kernels against globally bounded $E$, at cost
$O(N^{-K})$. On the retained shifts the induction hypothesis applies.
The mean value theorem, \eqref{eq:kernelderivative}, and
\eqref{eq:drift} give
\begin{equation}\label{eq:neighborbound}
 |((\mathcal A_{s_0}-\mathcal A_{s_y})E)(y)|
 \le C N^{-p}|s_0-s_y|+O(N^{-K})
 \le CbN^{-p-1}+O(N^{-K}).
\end{equation}
The derivative norm is uniform on the full interval between the two
saddles. There is no logarithmic loss in this estimate.
Therefore each retained $F(y)$ is bounded by
\[
 C\{N^{-K}+(1+b)N^{-p-1}\}.
\]
Insert this in \eqref{eq:frozenatbase}. The mass and first $b$ moment of
$\mathcal B_{s_0}$ are uniformly finite, while the discarded inverse tail
is $O(N^{-K})$. This gives
$|E(x)|\le C(N^{-K}+N^{-p-1})$, proving the lemma.
\end{proof}
The proof applies a translation-invariant inverse only after freezing
the parameter at the base state. It neither substitutes the wrong saddle
in a neighboring equation nor takes absolute values of an unproved
signed cancellation.

\subsection{Finite nested compact bootstrap}
Choose $K+1$ compact intervals
\[
 J_0\subset\operatorname{int}J_1\subset\cdots
     \subset\operatorname{int}J_K\subset(0,1),
\]
with the desired compact ratio set inside $J_0$ and the earlier residual
margin outside $J_K$. Use the same $H_K$ and the same bounded extension
throughout. Initially $E=O(1)$ on $J_K$. Apply
Lemma~\ref{lem:improvement} to obtain $O(N^{-1})$ on $J_{K-1}$,
then $O(N^{-2})$ on $J_{K-2}$, and so on, ending with
$E=O(N^{-K})$ on $J_0$. This proves Theorem~\ref{thm:allorders}.
There are only $K$ steps, independent of $N$; no scaled error bound has
been assumed in advance.

Finally, uniqueness as an asymptotic coefficient family follows by
comparing two continuous families successively along integer pairs whose
ratio tends to any prescribed interior value. At the first potentially
different order, multiply by the corresponding power of $N$ and pass
to the limit on a compact neighborhood. The two coefficient values
agree. This proves uniqueness on the whole interior, with no appeal to
a contraction subinterval.

\section{Diagonal normalization and the second correction}\label{sec:diagonal}
On the diagonal $N=2n$, $m=n$, the saddle is
$s_0=(1+\sqrt{17})/8$. From \eqref{eq:Q} and
Theorem~\ref{thm:allorders},
\[
 a_n=4n\frac{(4n-1)!!}{n!}\,
 \frac{f(t)h^nt^{-2n}}{\sqrt{2\pi nB}}\,\sqrt t
 \left(\sum_{j\ge0}D_j(0,0;s_0)n^{-j}\right)
 \left(\sum_{j\ge0}\frac{r_j(s_0)}{(2n)^j}\right),
\]
where each series is truncated at the desired fixed order with its proved
remainder. The factorial normalization is
\begin{equation}\label{eq:factorialnormal}
 \frac{(4n-1)!!}{(n!)^2}
 =\frac{16^n}{\pi\sqrt2\,n}\,\mathcal S(n^{-1}),\qquad
 \mathcal S(x)=1-\frac3{16}x+\frac9{512}x^2+\cdots.
\end{equation}
The logarithmic series of $\mathcal S$ has coefficient
\begin{equation}\label{eq:Stirlingseries}
 \frac{\mathsf B_{2v}}{2v(2v-1)}
 \left(4^{-(2v-1)}-2^{-(2v-1)}-2\right)
\end{equation}
at $x^{2v-1}$ and zero at positive even powers, where
$\mathsf B_{2v}$ denotes a Bernoulli number. Consequently
\begin{equation}\label{eq:diagonalalgorithm}
 \sum_{j\ge0}b_jx^j
 =\mathcal S(x)
  \left(\sum_{j\ge0}D_j(0,0;s_0)x^j\right)
  \left(\sum_{j\ge0}r_j(s_0)(x/2)^j\right).
\end{equation}
All its coefficients are in $\Q(s_0)=\Q(\sqrt{17})$.
The leading exponential and amplitude simplify to
\[
 d=\frac{16h}{t^2},\qquad
 C=\frac{2\sqrt t\,f(t)}{\pi^{3/2}\sqrt B},
\]
at $s=s_0$, which are precisely \eqref{eq:constants}.
Finite-order products of the established remainders prove
Corollary~\ref{cor:diagonal}.

For explicit second-order reproduction the algorithm yields
\begin{align}
 r_1(s)&=\frac{s^4+6s^3+7s^2+3s+4}
                    {4(s-1)(s+1)(s+2)^2},\label{eq:r1}\\
 D_1(0,0;s)&=\frac{s^6-7s^5-5s^4+32s^3+32s^2-9s-8}
                    {12(s-1)(s+1)(s+2)^3},\label{eq:D1}\\
 B_2(s)&=\frac{s^4+9s^3+13s^2+3s+4}{4t^2(s+2)^2}.
 \label{eq:B2}
\end{align}
To keep the higher numerators legible, write
\begin{align}
 P_2(s)={}&s^{10}+14s^9+22s^8-218s^7-1295s^6-3068s^5\notag\\
          &-3675s^4-2402s^3-1160s^2-500s-4,\label{eq:P2}\\
 T_2(s)={}&s^{12}-14s^{11}+759s^{10}+1574s^9-1079s^8-4386s^7\notag\\
          &-1346s^6+4410s^5+4128s^4-368s^3-1871s^2-576s+64.
 \label{eq:T2}
\end{align}
Then
\begin{equation}\label{eq:r2D2}
 r_2(s)=\frac{P_2(s)}{32s\,t^2(s+2)^5},\qquad
 D_2(0,0;s)=\frac{T_2(s)}{288t^2(s+2)^6}.
\end{equation}
These formulas can be checked directly from the finite algorithm without
inserting the proposed answers. One convenient Gaussian implementation
uses $\exp(\sum_{j\ge1}F_jz^j)=\sum_{n\ge0}E_nz^n$ with
\[
 F_j=\frac{g_j-r\kappa_j+\ell\mathbf1_{j=1}}{j!}(iu)^j
       +\frac{\kappa_{j+2}}{(j+2)!}(iu)^{j+2},\qquad
 E_0=1,\quad E_n=\frac1n\sum_{j=1}^n jF_jE_{n-j}.
\]
Apply \eqref{eq:Gaussianmoments} to $E_2,E_4$ and then divide by the
unshifted series. This independently generates $D_1,D_2$ and the ratio
coefficients needed in \eqref{eq:Aseries}.

For a direct check of the residual bookkeeping, set
$\delta=\alpha\ell-r$, $g=(\partial_\alpha\sqrt t)/\sqrt t$, and
\[
 \mathcal M(P)=\sqrt t\sum_{\ell\ge1,r}a_s(\ell,r)P(\ell,r).
\]
The known shifted terms, divided by the current $R_0$, are
\begin{align*}
 H_1&=r_1+g\delta,\\
 H_2&=\ell r_1+\delta(\partial_\alpha r_1+gr_1)
       +\ell g\delta+\frac{\delta^2}{2}(\partial_\alpha g+g^2).
\end{align*}
Thus the first two residual equations are
\begin{equation}\label{eq:explicitresiduals}
 r_1=-\mathcal M(A_1+g\delta)-B_1,\qquad
 r_2=-\mathcal M(H_2+A_1H_1+A_2)-B_2.
\end{equation}
This retains the derivatives of both $R_0$ and $r_1$, the change in
$1/(N-\ell)$, the cross term, and the opposite endpoint.

Finally, the second diagonal coefficient is
\begin{equation}\label{eq:b2product}
 b_2=\left[
 \frac9{512}+D_2+\frac{r_2}{4}-\frac3{16}D_1
       +\left(-\frac3{16}+D_1\right)\frac{r_1}{2}
 \right]_{s=s_0},
\end{equation}
which simplifies to \eqref{eq:b2}. Its decimal value is
$-0.583039199664459265278390709755\ldots$.
The saddle coefficients are in powers of $m^{-1}$, whereas $r_j$ are
in powers of $N^{-1}$; the factors $1/2$ and $1/4$ on the diagonal are
therefore necessary. The companion regenerates these identities exactly.

\section{All orders of the integer safe inverse}\label{sec:inverse}
Set
\begin{equation}\label{eq:Phi}
 c=\log(C\sqrt{2\pi}),\qquad
 \Phi(x)=x(\log(dx)-1)+c.
\end{equation}
Define $\lambda_j$ by the formal identity
\begin{equation}\label{eq:lambda}
 \sum_{j\ge1}\lambda_jx^{-j}
 =\log\left(\sum_{j\ge0}b_jx^{-j}\right)
       +\sum_{v\ge1}\frac{\mathsf B_{2v}}{2v(2v-1)}x^{-(2v-1)}.
\end{equation}
Each $\lambda_j\in\Q(\sqrt{17})$; specifically,
\begin{align}
 \lambda_1&=b_1+\frac1{12}
       =-\frac{2423}{6528}-\frac{907\sqrt{17}}{36992},\label{eq:lambda1}\\
 \lambda_2&=b_2-\frac{b_1^2}{2}
       =-\frac{4940681+601183\sqrt{17}}{10061824}.
 \label{eq:lambda2}
\end{align}
The forward theorem and Stirling's expansion imply, for every fixed $K$,
\begin{equation}\label{eq:loginteger}
 \log a_n=\Phi(n)+\sum_{j=1}^{K-1}\lambda_jn^{-j}+O(n^{-K}).
\end{equation}
The $\tfrac12\log n$ from $\log n!$ cancels the original
$n^{-1/2}$ factor.

For large $X$ put $L=\log X$, $T=L-c$, and
\begin{equation}\label{eq:x0}
 x_0=\frac{T}{W(dT/e)},\qquad
 \Psi_K(x)=\Phi(x)+\sum_{j=1}^{K-1}\lambda_jx^{-j},
\end{equation}
using the positive real Lambert $W$ branch. Thus $\Phi(x_0)=L$ and
$x_0\sim\log X/\log\log X$.

\begin{theorem}[All orders with the uncertainty inside the ceilings]
For every fixed $K\ge1$, the equation $\Psi_K(x)=L$ has a unique root
$x_K$ in $[x_0/2,2x_0]$ for sufficiently large $X$, and
$x_K-x_0=O((x_0\log x_0)^{-1})$. There is a constant $M_K>0$ such that
\begin{equation}\label{eq:ceilingbracket}
 \left\lceil x_K-\frac{M_K}{x_0^K\log(dx_0)}\right\rceil
 \le N(X):=\min\{n:a_n\ge X\}
 \le\left\lceil x_K+\frac{M_K}{x_0^K\log(dx_0)}\right\rceil
\end{equation}
for all sufficiently large $X$.
Starting at $y_0=x_0$, finitely many explicit Newton steps
\begin{equation}\label{eq:Newtoninverse}
 y_{v+1}=y_v-\frac{\Psi_K(y_v)-L}{\Psi_K'(y_v)}
\end{equation}
may replace $x_K$ in \eqref{eq:ceilingbracket}, after increasing $M_K$,
whenever $2^{v+1}-1\ge K$.
\end{theorem}
\begin{proof}
On $[x_0/2,2x_0]$, $\Psi_K'(x)\sim\log x_0$ uniformly, and its endpoints
bracket $L$ for large $X$. This proves existence and uniqueness there.
The value $\Psi_K(x_0)-L$ is $O(x_0^{-1})$, giving the claimed root
displacement by the mean value theorem.

No interpolation of the sequence is needed for integer rounding.
For integers in this interval, \eqref{eq:loginteger} has absolute error
at most $C_Kx_0^{-K}$. Moving a real argument by
$M_Kx_0^{-K}/\log(dx_0)$ changes $\Psi_K$ by more than this error when
$M_K$ is large enough. Every integer below the lower witness in
\eqref{eq:ceilingbracket} therefore has $a_n<X$, while the ceiling of the
upper witness has $a_n\ge X$. The forward equivalent gives
$a_{n+1}/a_n\sim dn$, so the sequence is eventually strictly increasing.
All earlier integers, including the finite initial prefix, are below
$X$ when $X$ is sufficiently large. These observations prove the bracket.

Throughout a neighborhood of $x_0$,
$|\Psi_K''/\Psi_K'|=O((x_0\log x_0)^{-1})$.
The initial error is $O((x_0\log x_0)^{-1})$, and the Newton error
recurrence therefore gives, for each fixed $v$,
\begin{equation}\label{eq:Newtonerror}
 |y_v-x_K|=O\left((x_0\log x_0)^{-(2^{v+1}-1)}\right).
\end{equation}
These iterates remain in the required neighborhood for sufficiently large
$X$. If $2^{v+1}-1\ge K$, their error is
$O(x_0^{-K}/\log x_0)$ and is absorbed by increasing the witness constant.
For example, $v=\lceil\log_2(K+1)\rceil$ always suffices.
\end{proof}
The constants and onset here are existence statements inherited from the
asymptotic remainder, not effective numerical certificates. A real error
tending to zero cannot be removed from the ceilings: arbitrarily close
approach to an integer prevents an unrounded $o(1)$ formula for the integer
threshold from following from these estimates.

\section{Reproducibility and scope}\label{sec:reproducibility}
The accompanying archive includes this complete LaTeX source, its PDF,
the unchanged Report 141 foundation, and an offline exact-arithmetic
companion. The companion regenerates the second correction, polynomial
boundary-zero identities, and frozen two-coordinate inverse checks from
computational definitions. It distinguishes typed exact fixtures from
public factual provenance. Its README specifies the runtime requirements,
isolated entrypoints, closed-inventory verification, absent external output
paths, semantic mutation tests, and deterministic PDF and ZIP workflows.
Finite computations do not establish the asymptotic remainder; the proof
is Sections~\ref{sec:saddle}--\ref{sec:fullrange}.

The crucial analytic points are the shifted saddle remainder uniform for
logarithmic shifts, summability of both endpoint expansions, the boundary
zeros in Lemma~\ref{lem:top}, the zero-index cancellation
\eqref{eq:Pzero}, and the base-frozen identity
\eqref{eq:frozenatbase} with the neighboring correction
\eqref{eq:neighborbound}. Together these give every fixed order on every
compact subset of the interior. They do not give convergence in the order,
boundary asymptotics, growing-order uniformity, or a certified numerical
threshold at a specified finite value of $X$.

\begin{thebibliography}{9}
\bibitem{foundation} \emph{A weighted Dyck path Newton diagonal}, Report 141,
October 2, 2026. Complete unchanged source and PDF included as the
foundation of this continuation. Contains the exact weighted-path bridge,
the formal convolution, source normalization for OEIS A258219, A258220,
and A292692, and the leading and first-corrected diagonal estimates.
\end{thebibliography}
\end{document}
